\documentclass[11pt]{article}

% Use Helvetica font.
%\usepackage{helvet}

%\usepackage[T1]{fontenc}
%\usepackage[sc]{mathpazo}
\usepackage{color}
\usepackage{amsmath,amsthm,amssymb,multirow,paralist}
%\renewcommand{\familydefault}{\sfdefault}

% Use 1/2-inch margins.
\usepackage[margin=0.8in]{geometry}
\usepackage{hyperref}

\begin{document}

\begin{center}
{\Large \textbf{Inductive Bias vs Global Attention: Comparing CNN and ViT architectures on TinyImageNet}

\vspace{10pt}

Project Midterm Report}

\vspace{10pt}

\textit{Team Members: Tyler Bibus}
\end{center}

% \linethickness{1mm}\line(1,0){498}

%%%%%%%%%%%%%%%%%%%%%%%%%%%%%%%%%%%%%%%%%%%%%%%%%%%%%%%%%%%%%%%%%%%%%%%%%%%%%%%

%%%%%%%%%%%%%%%%%%%%%%%%%%%%%%%%%%%%%%%%%%%%%%%%%%%%%%%%%%%%%%%%%%%%%%%%%%%%%%%

\begin{abstract}
  This project is on comparing Convolutional Neural Networks against Vision Transformers on TinyImageNet.
  The primary motivation behind this is to determine if similarly sized CNNs and ViTs perform dramatically differently.
  Convolutional neural networks used to dominate image classification on ImageNet, but now vision transformers are the new favorite.
  This project aims to see if CNNs still have a place against ViTs for image classification like on TinyImageNet.

  Currently TinyImageNet has been downloaded, normalized, and prepared with Pytorch and Pandas.
  An initial AlexNet-like model has been created with pytorch and has been trained with reasonable accuracies, with the best being 22.4\% validation accuracy.
\end{abstract}

\section{Introduction}

In this project we are using TinyImageNet, which is a downsampled subset of ImageNet.
This is ideal for our use case as we have limited compute budget at one GPU.

We plan on using two different archtiectures, that being an AlexNet-like Convolutional Neural Network and a Vision Transformer.
The AlexNet CNN is straight forward, being Convolutional layers with ReLU and MaxPooling followed by a few linear layers with dropouot.
The Vision transformer is more difficult to implement, as we have less direct experience with transformers. This transformer models will be modeled to have a similar amount of parameters as the CNN.

The expectation is that while AlexNet is still competitive, we expect the transformer to beat AlexNet. A more modern CNN like ResNet might be a slightly closer comparison as it has additional
functionality that enhances its performance.

\section{Related Work}

Please briefly introduce some related works to your project. The
related work could be some papers related to your machine
learning algorithm.

Usually, this part accounts for 0.5 page. Here, you can list one
or two related works for midterm. The final version could be
finished in the final report.

1. ImageNet Classification with Deep Convolutional Neural Networks.
This paper is the AlexNet paper which set CNNs as the gold standard for image classification pre-transformers.

2. An Image is Worth 16x16 Words: Transformers for Image Recognition at Scale
This paper is what brought transformers to image processing based on the successes transformers have had with Natural Language Processing.

\section{Your methods name}

In this section, you will introduce your machine learning
algorithms used in this part.

If you are using existing methods on a \textbf{new dataset}. You
can talk about ``dataset collection'', ``data pre-processing'',
``Feature selection'', ``Choices of hyper-parameter'',
``Evaluation metrics'', etc.

If you propose some \textbf{new methods}. You can talk about
``Drawbacks of existing methods'', ``Your proposed methods'',
``Your models based on your methods'', etc.

These are possible subsections in this part. You don't need to
include all of them. I want to see the contents that you spend
time on and have some insights.

This part is one of the most important sections. The length of this
section should be about 1.5 pages.

\section{Preliminary Results}

In this section, you will list some experimental results with
some comparisons with some baseline methods.

For example, you choose to use SVM on a new dataset. Possible
baselines would be linear perceptron and logistic regression by
using the same set of features.

In this midterm report, you can use a table to show some
preliminary results along with some discussions. This part should
be about 0.5 page.

\section{Future plan}

In this section, briefly discuss your plan based on current
achievements.

\section{References}

Please include your references here.


\end{document}
